%=====================================================================
\chapter{Knowledge Representation}\label{ch:kr}
%=====================================================================
\locator{4}{Part IV --- Knowledge, Logic and Reasoning}

\begin{outcomes}
By the end of this chapter you will be able to:
\begin{tight}
  \item \textbf{State} what a knowledge representation must support, and
        \textbf{judge} a candidate representation against those criteria.
  \item \textbf{Build} a semantic network and \textbf{answer} queries by
        inheritance.
  \item \textbf{Distinguish} frames, slots and defaults, and \textbf{explain} what
        default override buys and what it costs.
  \item \textbf{Identify} a multiple-inheritance conflict and \textbf{say} why no
        purely local rule resolves it.
  \item \textbf{Describe} what an ontology adds, and \textbf{say} where this course
        hands over to \emph{Knowledge Representation and Reasoning}.
\end{tight}
\end{outcomes}

\begin{prereq}
\chref{python} (dictionaries, sets and recursion) and \chref{introduction} (the
third test of the checklist: an explicit representation a domain expert could read
and correct). No logic is assumed --- \chref{propositional} makes formal what this
chapter introduces informally.
\end{prereq}

\begin{keyterms}
knowledge representation $\bullet$ ontological commitment $\bullet$
epistemological adequacy $\bullet$ semantic network $\bullet$ node and link
$\bullet$ \texttt{is-a} $\bullet$ \texttt{instance-of} $\bullet$ inheritance
$\bullet$ frame $\bullet$ slot $\bullet$ filler $\bullet$ default $\bullet$
override $\bullet$ non-monotonic reasoning $\bullet$ multiple inheritance
$\bullet$ ontology $\bullet$ taxonomy
\end{keyterms}

\section{The Change of Question}

Parts~II and~III asked \emph{what should I do}. Every algorithm took a goal and
searched for a sequence of actions. The knowledge was implicit: buried in a
successor function, a constraint predicate, an operator table.

Part~IV asks a different question: \emph{what follows from what I have been told}.
The delivery team knows a great many things --- that a hotfix is a kind of release,
that every release needs a sign-off, that Ayesha owns the payments service, that
releases normally go out on Thursdays except this one. None of that is a goal. It
is what the organisation knows, and the useful question is what else must therefore
be true.

That change is why this part exists, and it has a practical consequence worth
stating at the start: \emph{a representation can be corrected by someone who cannot
program}. An operations lead can tell you that a hotfix does not need a full
regression pass. If that fact lives in a data structure, they can point at it. If it
lives in an \texttt{if} statement, they cannot.

\begin{definitionbox}[What a knowledge representation must support]
A representation is a scheme for writing down what is known so that a program can
reason with it. Four requirements, from Russell and Norvig:
\begin{tight}
  \item \term{Representational adequacy} --- it can express the things you need to
        say, including the awkward ones: defaults, exceptions, uncertainty.
  \item \term{Inferential adequacy} --- new facts can be derived from stored ones.
  \item \term{Inferential efficiency} --- the derivation happens in acceptable
        time.
  \item \term{Acquisitional efficiency} --- new knowledge can be added easily, and
        ideally by a domain expert rather than by an engineer.
\end{tight}
These pull against each other. A representation expressive enough to say everything
is generally one in which inference is intractable, and that tension is the whole
subject.
\end{definitionbox}

\section{Semantic Networks}

The oldest and simplest representation: a graph. Nodes are concepts or individuals;
labelled edges are relations between them. Two link types carry most of the weight:
\texttt{is-a}, relating a class to a more general class, and \texttt{instance-of},
relating an individual to its class.

\dsfig{%
\begin{tikzpicture}[font=\scriptsize,
  cls/.style={rounded corners=3pt, draw=primary, line width=1pt, fill=primary!10,
              text=primary!60!black, font=\scriptsize\bfseries, align=center,
              minimum width=2.1cm, minimum height=0.6cm},
  ind/.style={rounded corners=10pt, draw=goldhl, line width=1pt, fill=goldhl!12,
              text=goldhl!50!black, font=\scriptsize\bfseries, align=center,
              minimum width=2.1cm, minimum height=0.6cm},
  isa/.style={-{Stealth[length=1.9mm]}, primary!60, line width=0.9pt},
  inst/.style={-{Stealth[length=1.9mm]}, goldhl!70, line width=0.9pt,
               dash pattern=on 3pt off 2pt},
  rel/.style={-{Stealth[length=1.9mm]}, tealhl!70, line width=0.9pt},
  el/.style={font=\tiny, text=neutral, inner sep=1.5pt}]

\node[cls] (art)     at (0,3.0)    {Artefact};
\node[cls] (rel)     at (-2.4,1.8) {Release};
\node[cls] (doc)     at (2.4,1.8)  {Document};
\node[cls] (hot)     at (-4.2,0.6) {Hotfix};
\node[cls] (maj)     at (-0.9,0.6) {Major Release};
\node[ind] (r47)     at (-4.2,-0.7) {Release 4.7};
\node[cls] (dev)     at (3.6,0.4)  {Developer};
\node[ind] (ayesha)  at (3.6,-0.7) {Ayesha};
\node[cls] (svc)     at (0.9,-0.7) {Service};
\node[ind] (pay)     at (0.9,-1.9) {Payments};

\draw[isa] (rel) -- node[el, above left] {is-a} (art);
\draw[isa] (doc) -- node[el, above right] {is-a} (art);
\draw[isa] (hot) -- node[el, left] {is-a} (rel);
\draw[isa] (maj) -- node[el, right] {is-a} (rel);
\draw[inst] (r47) -- node[el, left] {instance-of} (hot);
\draw[inst] (ayesha) -- node[el, right] {instance-of} (dev);
\draw[inst] (pay) -- node[el, left] {instance-of} (svc);
\draw[rel] (ayesha) -- node[el, above] {owns} (svc);
\draw[rel] (r47) to[bend right=18] node[el, below] {changes} (pay);

\node[rounded corners=3pt, fill=lightbg, draw=primary!30, line width=0.8pt,
      text=primary!70!black, font=\scriptsize, text width=5.4cm,
      align=flush left, inner sep=6pt, anchor=west] at (5.2,2.1)
  {\textbf{The query that matters.} Nothing in this graph says
   \emph{Release 4.7 is an Artefact}. It is derived, by walking
   \texttt{instance-of} once and then \texttt{is-a} upwards until the walk stops:
   Release 4.7 $\rightarrow$ Hotfix $\rightarrow$ Release $\rightarrow$ Artefact.
   That walk is \term{inheritance}, and it is the entire inference mechanism of a
   semantic network.};
\end{tikzpicture}}%
{A fragment of a delivery ontology. Rounded boxes are individuals, square boxes are
classes; solid arrows are \texttt{is-a}, dashed are \texttt{instance-of}.}
{semantic-net}

Inheritance is cheap, and it is the whole reason the representation is worth
having: a property stated once at \texttt{Release} is available to every hotfix and
every major release, without being repeated and --- more importantly --- without
being repeated \emph{inconsistently}.

\section{Frames, Slots and Defaults}

A semantic network drawn as a graph becomes unreadable at a few hundred nodes. The
same information organised by concept is a \term{frame}: a named bundle of
\term{slots}, each with a \term{filler} or a \term{default}.

\begin{examplebox}[Three frames from the delivery domain]
\begin{tabular}{@{}L{3.0cm}L{5.1cm}L{5.9cm}@{}}
\toprule
\textbf{Frame} & \textbf{Slot} & \textbf{Value} \\
\midrule
\texttt{Release} & \texttt{is-a} & \texttt{Artefact} \\
 & \texttt{needs-signoff} & \texttt{yes} \emph{(default)} \\
 & \texttt{regression-pass} & \texttt{full} \emph{(default)} \\
 & \texttt{window} & \texttt{Thursday} \emph{(default)} \\
\midrule
\texttt{Hotfix} & \texttt{is-a} & \texttt{Release} \\
 & \texttt{regression-pass} & \texttt{smoke} \quad\textbf{(overrides)} \\
 & \texttt{window} & \texttt{any} \quad\textbf{(overrides)} \\
\midrule
\texttt{Release 4.7} & \texttt{instance-of} & \texttt{Hotfix} \\
 & \texttt{changes} & \texttt{Payments} \\
\bottomrule
\end{tabular}
\end{examplebox}

Ask for \texttt{Release 4.7}'s \texttt{regression-pass}. It has none of its own, so
the search goes to \texttt{Hotfix}, which says \texttt{smoke}, and stops. Ask for
\texttt{needs-signoff}: \texttt{Release 4.7} has none, \texttt{Hotfix} has none,
\texttt{Release} says \texttt{yes}. The rule is \emph{the most specific value
wins}, and it falls out of searching upward and stopping at the first hit.

\begin{conceptbox}
Defaults are what make this representation fit an organisation, and they are also
where it stops being logic.

``Releases go out on Thursdays'' is not a universally quantified statement --- it
is false of hotfixes --- and yet it is exactly what the delivery lead would tell
you, and exactly what you want the system to assume when it knows nothing more.
Reasoning that can be \emph{withdrawn} when a more specific fact arrives is
\term{non-monotonic}: adding knowledge can remove a conclusion.

Classical logic, the subject of \chref{propositional} and \chref{fol}, is
monotonic: once entailed, always entailed. The two chapters therefore buy soundness
and lose the ability to say ``normally''. That trade is real, and neither side wins
outright.
\end{conceptbox}

\section{When Inheritance Has No Answer}

Defaults are convenient until a concept inherits from two places that disagree.

\dsfig{%
\begin{tikzpicture}[font=\scriptsize,
  cls/.style={rounded corners=3pt, draw=#1, line width=1pt, fill=#1!10,
              text=#1!55!black, font=\scriptsize\bfseries, align=center,
              minimum width=2.6cm, minimum height=0.68cm},
  ind/.style={rounded corners=10pt, draw=goldhl, line width=1.2pt, fill=goldhl!14,
              text=goldhl!50!black, font=\scriptsize\bfseries, align=center,
              minimum width=2.6cm, minimum height=0.68cm},
  isa/.style={-{Stealth[length=1.9mm]}, line width=1pt},
  el/.style={font=\tiny, text=neutral, inner sep=1.5pt}]

\node[cls=primary]  (rel) at (0,2.1)     {Release};
\node[cls=greenhl]  (hot) at (-2.3,0.75) {Hotfix};
\node[cls=accent]   (reg) at (2.3,0.75)  {Regulated Change};
\node[ind]          (r47) at (0,-0.6)    {Release 4.7};

\draw[isa, primary!60] (hot) -- (rel);
\draw[isa, primary!60] (reg) -- (rel);
\draw[isa, greenhl!70] (r47) -- node[el, left] {is-a} (hot);
\draw[isa, accent!70]  (r47) -- node[el, right] {is-a} (reg);

\node[font=\tiny, text=greenhl!50!black, anchor=east, align=right]
  at (-3.7,0.75) {\texttt{needs-signoff}\\ \textbf{no}};
\node[font=\tiny, text=accent!55!black, anchor=west, align=left]
  at (3.75,0.75) {\texttt{needs-signoff}\\ \textbf{yes}};
\node[font=\tiny, text=primary!60!black, anchor=west, align=left]
  at (1.5,2.1) {\texttt{needs-signoff}: yes};

\node[rounded corners=3pt, fill=accent!7, draw=accent!55, line width=0.9pt,
      text=accent!55!black, font=\scriptsize, text width=12.8cm,
      align=flush left, inner sep=6pt] at (0,-2.1)
  {\textbf{Does Release 4.7 need a sign-off?} Both parents are equally specific and
   they disagree. Depth-first inheritance answers whichever parent it happens to
   visit first --- so the answer depends on the order the links were entered, which
   is not a property of the knowledge at all. Breadth-first inheritance has the
   same defect with a different bias.

   \smallskip
   There is no purely local rule that settles this. The honest options are to
   \textbf{detect the conflict and refuse to answer}, to demand an explicit
   precedence between the two parents, or to move to a representation with a
   defined semantics for it. A system that quietly returns ``no'' here is the
   dangerous one --- and in this instance it would be waving a regulated change
   through without a sign-off.};
\end{tikzpicture}}%
{A multiple-inheritance conflict --- the ``Nixon diamond'' of the literature, in a
delivery setting. The lab detects it rather than guessing.}{diamond}

\section{Ontologies}

An \term{ontology} is a semantic network taken seriously: a shared, formal
vocabulary of the types in a domain, their relations, and the constraints on them,
agreed across an organisation rather than invented per system.

The value is not the graph; it is the \emph{agreement}. When the delivery tool, the
incident tool and the finance system all mean the same thing by ``release'', a query
can cross them. When they do not, the integration work is endless and the reports
disagree. Most of the practical benefit of ontologies in industry comes from that,
and not from the inference.

Formally, an ontology commits to answers for questions a semantic network leaves
open: is \texttt{is-a} transitive (yes), can an individual belong to two classes
(usually yes), can a class be an individual, what makes two individuals identical.
Those are \term{ontological commitments}, and writing them down is most of the work.

\section{Where This Course Stops}

This chapter is a first pass, deliberately. Everything after it in this handout
uses semantic networks and frames as the \emph{thing reasoned over}:
\chref{propositional} and \chref{fol} give inference a formal footing, and
\chref{expert} turns a frame base plus rules into a working system.

What this course does not cover is the formal side of representation itself:
description logic, reasoning over description logic, ontology languages such as OWL,
logical agents and commonsense reasoning. Those are the content of the HEC elective
\emph{Knowledge Representation and Reasoning}, for which this course is the
prerequisite; Appendix~F sets out the line.

\begin{worked}[Answering four queries by inheritance]
Using the frames of the examplebox and the conflict of \dsref{diamond}.

\smallskip
\textbf{Query 1: \texttt{Release 4.7 . regression-pass}}
\begin{tight}
  \item \texttt{Release 4.7}: no local slot. Follow \texttt{instance-of} to
        \texttt{Hotfix}.
  \item \texttt{Hotfix}: \texttt{regression-pass = smoke}. \textbf{Stop.}
\end{tight}
\textbf{Answer: \texttt{smoke}}, inherited from \texttt{Hotfix}. Note that
\texttt{Release} says \texttt{full} and is never consulted --- the search stops at
the first hit, which is what ``most specific wins'' means operationally.

\smallskip
\textbf{Query 2: \texttt{Release 4.7 . needs-signoff}} (ignoring the second parent
for the moment)
\begin{tight}
  \item \texttt{Release 4.7}: nothing. \texttt{Hotfix}: nothing.
        \texttt{Release}: \texttt{yes}. \textbf{Stop.}
\end{tight}
\textbf{Answer: \texttt{yes}}, inherited from two levels up. The same walk answered
both queries; only the depth at which it stopped differed.

\smallskip
\textbf{Query 3: \texttt{Release 4.7 . window}}
\begin{tight}
  \item \texttt{Hotfix} overrides with \texttt{any}. \textbf{Answer: \texttt{any}}
        --- the default \texttt{Thursday} at \texttt{Release} is correctly
        suppressed.
\end{tight}
Now add the fact that \texttt{Release 4.7} is also a \texttt{Regulated Change} whose
window is \texttt{Thursday}. The answer changes to a conflict. \emph{Adding
knowledge removed a conclusion}: that is non-monotonicity, in three lines.

\smallskip
\textbf{Query 4: \texttt{Release 4.7 . needs-signoff}, with both parents.}
\begin{tight}
  \item \texttt{Hotfix} says \texttt{no}; \texttt{Regulated Change} says
        \texttt{yes}. Both are at distance~1 from the individual.
  \item Depth-first inheritance returns whichever was entered first. That is not an
        answer; it is an artefact of the file.
\end{tight}
\textbf{Answer: conflict}, and the system should say so. The lab returns the set
$\{\texttt{yes}, \texttt{no}\}$ with its provenance, and leaves the decision to a
person --- which in this case is the only defensible behaviour, because the two
sources are a delivery convention and a regulatory obligation and they are not
equally negotiable.
\end{worked}

\begin{pitfall}
\begin{tight}
  \item \textbf{Letting inheritance guess on a multiple-inheritance conflict.} The
        answer then depends on insertion order. Detect it and refuse.
  \item \textbf{Treating a default as a universal.} ``Releases go out on Thursdays''
        is false of hotfixes. If you write it as $\forall$ in \chref{fol}, the
        knowledge base becomes inconsistent the first time a hotfix appears.
  \item \textbf{Putting knowledge in code.} A rule inside an \texttt{if} statement
        cannot be reviewed by the person who knows whether it is right.
  \item \textbf{Cycles in \texttt{is-a}.} An inheritance walk that does not keep a
        visited set will not terminate. The lab keeps one.
  \item \textbf{Confusing \texttt{is-a} with \texttt{instance-of}.} A hotfix
        \emph{is a kind of} release; Release 4.7 \emph{is a} hotfix. Conflating
        them lets a class inherit properties of an individual, and the errors are
        very hard to see.
  \item \textbf{Expecting the ontology to pay for itself through inference.} In
        practice most of the value is the shared vocabulary.
\end{tight}
\end{pitfall}

%---------------------------------------------------------------------
\begin{lab}[A Semantic Network That Answers Queries]
Frames in a dictionary, inheritance by walking upward, defaults with override, and
a conflict that is detected rather than guessed.
\begin{lstlisting}[language=Python]
"""Chapter 10 lab -- a semantic network with frames and defaults.

Inheritance is one upward walk. The interesting cases are the third
query, where a default is correctly suppressed by a more specific one,
and the fourth, where two equally specific parents disagree and the
only defensible answer is to say so.
"""

# frame -> {slot: value}. 'is-a' and 'instance-of' are the links the
# inheritance walk follows; everything else is ordinary knowledge.
FRAMES = {
    "Artefact": {
        "owned-by": "Delivery",
    },
    "Release": {
        "is-a": ["Artefact"],
        "needs-signoff": "yes",          # default
        "regression-pass": "full",       # default
        "window": "Thursday",            # default
    },
    "Hotfix": {
        "is-a": ["Release"],
        "regression-pass": "smoke",      # overrides Release
        "window": "any",                 # overrides Release
    },
    "Regulated Change": {
        "is-a": ["Release"],
        "needs-signoff": "yes",
        "window": "Thursday",
    },
    "Release 4.7": {
        "instance-of": ["Hotfix"],
        "changes": "Payments",
    },
    "Developer": {},
    "Ayesha": {"instance-of": ["Developer"], "owns": "Payments"},
}

LINKS = ("instance-of", "is-a")


def parents(frame):
    out = []
    for link in LINKS:
        out += FRAMES.get(frame, {}).get(link, [])
    return out


def ancestors(frame):
    """Every class above this one, nearest first. The seen set is what
    stops an is-a cycle from looping forever."""
    out, seen, queue = [], {frame}, list(parents(frame))
    while queue:
        f = queue.pop(0)
        if f in seen:
            continue
        seen.add(f)
        out.append(f)
        queue += parents(f)
    return out


def get(frame, slot):
    """Most specific value wins. Returns (value, where it came from),
    or a conflict if two equally distant sources disagree."""
    if slot in FRAMES.get(frame, {}):
        return FRAMES[frame][slot], frame

    # breadth-first by distance, so 'equally specific' is well defined
    level, seen = list(parents(frame)), {frame}
    distance = 1
    while level:
        hits = []
        for f in level:
            if f in seen:
                continue
            if slot in FRAMES.get(f, {}):
                hits.append((FRAMES[f][slot], f))
        values = {v for v, _ in hits}
        if len(values) > 1:
            return ("CONFLICT", sorted(hits)), distance
        if hits:
            return hits[0][0], hits[0][1]
        seen |= set(level)
        nxt = []
        for f in level:
            nxt += parents(f)
        level = [f for f in nxt if f not in seen]
        distance += 1
    return None, None


def isa_chain(frame):
    return [frame] + ancestors(frame)


if __name__ == "__main__":
    print("inheritance chain for Release 4.7:")
    print("   " + " -> ".join(isa_chain("Release 4.7")))
    assert "Artefact" in ancestors("Release 4.7")
    print()

    print("queries, with the frame each answer came from")
    print("-" * 56)
    for slot in ("regression-pass", "needs-signoff", "window",
                 "changes", "owned-by"):
        val, src = get("Release 4.7", slot)
        print("   %-16s = %-10s  (from %s)" % (slot, val, src))

    # the default was correctly suppressed
    assert get("Release 4.7", "regression-pass") == ("smoke", "Hotfix")
    assert get("Release", "regression-pass") == ("full", "Release")
    # and a default two levels up is still found
    assert get("Release 4.7", "needs-signoff") == ("yes", "Release")
    print()
    print("'smoke' overrode 'full'; 'yes' was still reached two levels up.")

    # --- now add a second, equally specific parent ----------------
    print()
    print("adding 'Release 4.7 is-a Regulated Change' ...")
    FRAMES["Hotfix"]["needs-signoff"] = "no"      # delivery convention
    FRAMES["Release 4.7"]["is-a"] = ["Regulated Change"]

    val, dist = get("Release 4.7", "needs-signoff")
    print("   needs-signoff =", val[0] if isinstance(val, tuple) else val)
    assert isinstance(val, tuple) and val[0] == "CONFLICT"
    for v, src in val[1]:
        print("      %-4s from %s" % (v, src))
    print()
    print("Two parents at the same distance disagree. Depth-first")
    print("inheritance would answer whichever link was entered first,")
    print("which is a property of the file and not of the knowledge.")
    print("Here that would wave a regulated change through with no")
    print("sign-off, so the only defensible answer is to report the")
    print("conflict and its provenance.")

    # non-monotonicity, demonstrated: adding knowledge REMOVED an answer
    print()
    print("Note what just happened: adding a fact removed a")
    print("conclusion. Classical logic cannot do that, which is")
    print("exactly the trade Chapters 11 and 12 make.")
\end{lstlisting}
Then add an explicit precedence --- a slot \texttt{overrides} on
\texttt{Regulated Change} naming \texttt{Hotfix} --- and make \texttt{get} consult
it before reporting a conflict. That is the smallest honest fix, and it is worth
noticing that it works by adding \emph{knowledge about the knowledge}, which is the
direction description logic goes in.
\end{lab}

%---------------------------------------------------------------------
\begin{checkpoint}
\begin{tightnum}
  \item \texttt{Release} says \texttt{regression-pass = full} and \texttt{Hotfix}
        says \texttt{smoke}. Explain, in terms of the inheritance walk, why
        \texttt{Release 4.7} gets \texttt{smoke}, and why no rule about
        ``specificity'' has to be stated separately.
  \item Adding the fact that Release~4.7 is a regulated change \emph{removed} the
        system's answer about sign-off. Name the property of this kind of reasoning
        that this demonstrates, and say whether classical logic has it.
  \item Give one thing a semantic network can express that a set of universally
        quantified logical statements cannot, and one thing logic can do that the
        network cannot.
\end{tightnum}
\end{checkpoint}

%---------------------------------------------------------------------
\begin{chaptersummary}
\begin{tight}
  \item A knowledge representation must be \term{representationally} and
        \term{inferentially adequate}, \term{efficient}, and easy to add to --- and
        those requirements pull against each other. Expressive representations tend
        to have intractable inference.
  \item A \term{semantic network} is nodes and labelled links. \texttt{is-a}
        relates a class to a more general class, \texttt{instance-of} relates an
        individual to a class, and \term{inheritance} --- one upward walk --- is
        the whole inference mechanism.
  \item A \term{frame} organises the same knowledge by concept: slots with fillers
        or \term{defaults}. The most specific value wins, which falls out of
        stopping the walk at the first hit.
  \item Defaults make the representation fit an organisation and make the reasoning
        \term{non-monotonic}: adding knowledge can remove a conclusion. Classical
        logic is monotonic and cannot say ``normally''.
  \item \term{Multiple inheritance} from disagreeing parents has no purely local
        resolution. Detect the conflict, report its provenance, and refuse to
        guess --- the dangerous system is the one that quietly answers.
  \item An \term{ontology} is a shared, formal vocabulary with explicit
        \term{ontological commitments}. In practice most of its value is the
        agreement, not the inference.
  \item This is a first pass. Description logic, ontology languages and commonsense
        reasoning belong to \emph{Knowledge Representation and Reasoning}
        (Appendix~F).
\end{tight}
\end{chaptersummary}

\begin{reviewq}
  \item \textbf{(MCQ)} In a semantic network, \texttt{is-a} relates:
        (a)~an individual to a class (b)~a class to a more general class
        (c)~two individuals (d)~a slot to a filler.
  \item \textbf{(MCQ)} ``The most specific value wins'' is implemented by:
        (a)~sorting all values (b)~stopping the upward walk at the first hit
        (c)~preferring the root (d)~taking the majority.
  \item \textbf{(MCQ)} Reasoning in which adding a fact can remove a conclusion is:
        (a)~monotonic (b)~non-monotonic (c)~unsound (d)~incomplete.
  \item \textbf{(MCQ)} A multiple-inheritance conflict between two equally specific
        parents should be: (a)~resolved depth-first (b)~resolved breadth-first
        (c)~reported, with its provenance (d)~resolved by the alphabetically first
        parent.
  \item \textbf{(Short)} State the four requirements on a knowledge representation
        and give one pair that pull against each other, with an example.
  \item \textbf{(Short)} Explain the difference between \texttt{is-a} and
        \texttt{instance-of}, and describe a concrete error caused by conflating
        them.
  \item \textbf{(Short)} Why is ``most of the value of an ontology is the
        agreement, not the inference'' a reasonable claim? Give a delivery example.
  \item \textbf{(Applied)} Extend the frames of this chapter with a class
        \texttt{Emergency Hotfix}, a kind of \texttt{Hotfix}, whose
        \texttt{needs-signoff} is \texttt{retrospective}. Give the value of
        \texttt{needs-signoff}, \texttt{regression-pass} and \texttt{window} for an
        instance of it, naming the frame each answer comes from. Then say what
        happens if that instance is also a \texttt{Regulated Change}, and what the
        system should do.
\end{reviewq}
